\begin{proposition}[E.26 --- Proved: the exact kernel in the once-punctured comparison]
\label{E-elliptic-cover-comparison}
Let $P$ be the geometric pro-three group of $E\setminus\{O\}$,
with free generators $x,y$, and let
$N=\ker(P\to(\mathbb Z/3)^2)$. The covering
$[3]:E\setminus E[3]\to E\setminus\{O\}$ identifies
this $N$ with the geometric group used in E.18--E.25.
The natural comparison of its class-two quotient with
the metabelian quotient of $P$ is the diagram
\[
 N/\gamma_3N\ \longrightarrow\
 N/(P''\gamma_3N)\ \longleftarrow\ N/P''
 \ \lhook\joinrel\longrightarrow\ P/P'' .
 \tag{E26.1}
\]
The first arrow has a nonzero kernel, detected in a
Heisenberg group of order $27$. In particular, the
quotient data on the right of (E26.1) do not specify
all the elements of the quotient on the left through
these comparison maps.
\end{proposition}
\begin{proof}
The degree-nine isogeny is finite \emph{\'etale} on the
displayed punctured curves. Its deck group is $E[3]$,
and the monodromy is the compact homology modulo three.
For the once-punctured elliptic group its abelianization
is precisely compact homology. This proves the kernel
description and specifies the covering identification.
Since $P'\subset N$, one has $P''\subset N$ and all
the maps in (E26.1) are the stated continuous quotient
or inclusion maps.

Put $c=xyx^{-1}y^{-1}$ and
$z=[c,xcx^{-1}]$. Both entries in this commutator
belong to $P'$, so $z\in P''$.
For an explicit free basis of $N$, use the cover of the
two-petal graph with vertices $(i,j)\in\mathbb F_3^2$.
Its positive $x$-edges increment $i$ and its positive
$y$-edges increment $j$. Choose as spanning tree the
two edges $(i,0)\to(i+1,0)$ for $i=0,1$ and the six
edges $(i,j)\to(i,j+1)$ for $i=0,1,2$, $j=0,1$.
The ten remaining oriented edges give a free basis
$e_0,\ldots,e_9$ after the tree is contracted.
Order them lexicographically by $(i,j,a)$, with
$a=1$ for $x$ and $a=2$ for $y$. Directly tracing the
two closed words gives
\[
 [c]_{N^{\rm ab}}=-e_0,\qquad
 [xcx^{-1}]_{N^{\rm ab}}=-e_3.
\]
Their exterior product is the primitive vector
$e_0\wedge e_3$. Hence $z$ is nonzero in
$\gamma_2N/\gamma_3N$.

The finite certificate sends $e_0$ and $e_3$ to
$(1,0,0)$ and $(0,1,0)$ in the group on
$\mathbb F_3^3$ with multiplication
\[
 (a,b,d)(a',b',d')=(a+a',b+b',d+d'+ab'),
\]
and sends every other basis generator to the identity.
The image of $z$ is $(0,0,1)$. The executed
\texttt{elliptic\_cover\_kernel.py} reconstructs the
Schreier words and verifies this exact value. Since
the finite target has class two, it also proves that
$z\notin\gamma_3N$ profinitely. Its image in the
middle quotient of (E26.1) is zero because $z\in P''$.
This proves the asserted nonzero kernel.

This comparison is the precise obstacle encountered
in trying to evaluate the remaining class of E.25
directly from the once-punctured elliptic power series
of Nakamura, Theorem A (4.12), and the path-torsor
formulas of Wojtkowiak~\cite{Nakamura1995Elliptic,Wojtkowiak2004Paths}.
Those primary formulas were opened and their domains
checked: the former is a Galois character construction
over the field of all prime-power division points,
and the latter begins with paths on
$E\setminus\{O\}$. Applying them to the actual
nine-puncture extension requires an additional
compatible construction retaining the kernel just
exhibited. The finite example does not assert that
no further comparison using these formulas exists;
it proves exactly what the displayed quotient maps
retain and remove. The anchor is \S51.16, folios
733v--736v, scans 172--178, and LM7-E.23(i).
The supplied LM7 tower return, E.50, gives a larger
kernel calculation for this same cover. The argument
above is an independent finite certificate for its
degree-two part; E.30 verifies its degree-three kernel
and locates the remaining arithmetic class in that kernel.
\end{proof}
